\section{Deviations from the frozen protocol}
\label{app:deviations}

\begin{enumerate}
\item \textbf{Download date (2026-10-28).} The post-release Tier B and W data were downloaded on the evening of the freeze day, from
about 18:10 UTC, instead of ``on or after 2026-10-29'', because the author was unavailable the next day and the GPU runs had to
start. The freeze commit (17:59:59 UTC) precedes the download; the test windows end on 2026-10-20 and were fully covered by ERA5;
nothing else changed.
\item \textbf{Variate-attention ablation (found 2026-10-30).} \texttt{tfm3-novarattn} set the flag
\texttt{use\_variate\_attention = False} on the top module of the loaded model, where it is only stored; the forward pass reads it on
each of the 20 layers. Its forecasts therefore equalled those of \texttt{tfm3} up to GPU non-determinism. The error was found while
drafting this paper, after the confirmatory analysis, because the ablation's ratio was exactly 1 on cells where multivariate mode
changes the error by up to 55\%. The other module ablations read their flags at run time and are valid. The override now sets the
flag on every submodule that has it, the smoke test checks that each switch changes a multivariate forecast, and the ablation is
re-run on the same tasks with the corrected code; family S5 is recomputed with it. Families P and H do not involve this
configuration.
\end{enumerate}

\section{Tier A per task}
\label{app:tiera}

\begin{table}[h]
\centering
\caption{Tier A, \texttt{tfm3-eval}: our SQL and Google's published SQL per task, relative deviations of SQL and MASE, pipeline
controls (Seasonal Naive and Chronos-2), and SQL relative to fev's published statistical ensemble.}
\scriptsize
\setlength{\tabcolsep}{3pt}
\begin{tabular}{lrrrrlr}
\toprule
Task & SQL ours & SQL Google & $\Delta$SQL & $\Delta$MASE & Controls & SQL / stat.\ ens. \\
\midrule
ETT\_15T & 0.543 & 0.543 & +0.00\% & +0.15\% & yes & -- \\
LOOP\_SEATTLE\_1D & 0.759 & 0.759 & +0.00\% & +0.00\% & yes & 0.925 \\
SZ\_TAXI\_15T & 0.392 & 0.392 & +0.00\% & -0.00\% & yes & -- \\
SZ\_TAXI\_1H & 0.401 & 0.401 & +0.00\% & +0.00\% & no & 0.581 \\
australian\_tourism & 0.791 & 0.791 & +0.00\% & -0.00\% & no & 1.083 \\
boomlet\_1062 & 0.545 & 0.545 & +0.04\% & -0.01\% & yes & 0.554 \\
boomlet\_1487 & 0.412 & 0.406 & +1.42\% & +1.27\% & yes & 0.605 \\
boomlet\_2187 & 0.714 & 0.714 & +0.00\% & +0.14\% & yes & 0.561 \\
ecdc\_ili & 2.145 & 2.145 & -0.00\% & -0.01\% & no & 0.559 \\
entsoe\_1H & 0.315 & 0.316 & -0.24\% & -0.51\% & no & 0.354 \\
epf\_pjm & 0.411 & 0.411 & -0.08\% & +0.02\% & yes & 0.844 \\
ercot\_1D & 0.780 & 0.780 & +0.00\% & -0.00\% & no & 0.621 \\
ercot\_1M & 0.764 & 0.764 & -0.00\% & -0.01\% & yes & 1.003 \\
ercot\_1W & 0.911 & 0.911 & -0.00\% & -0.00\% & no & 0.435 \\
favorita\_stores\_1M & 1.923 & 1.813 & +6.11\% & +4.92\% & yes & 0.990 \\
fred\_md\_2025/macro & 5.110 & 5.080 & +0.60\% & +0.48\% & no & 0.890 \\
gvar & 0.568 & 0.567 & +0.19\% & +0.19\% & no & 0.964 \\
hospital & 0.682 & 0.682 & +0.00\% & +0.00\% & no & 0.979 \\
jena\_weather\_10T & 0.338 & 0.340 & -0.45\% & -0.25\% & yes & 0.610 \\
kdd\_cup\_2022\_10T & 0.425 & 0.429 & -0.99\% & -0.41\% & no & 0.593 \\
m5\_1W & 0.878 & 0.878 & -0.03\% & -0.05\% & no & 0.938 \\
proenfo\_gfc17 & 0.403 & 0.407 & -1.02\% & -0.99\% & no & 0.353 \\
restaurant & 0.700 & 0.700 & -0.00\% & +0.00\% & no & 0.988 \\
rohlik\_sales\_1W & 1.114 & 1.215 & -8.32\% & -9.09\% & no & 0.677 \\
rossmann\_1W & 0.265 & 0.297 & -10.61\% & -9.10\% & no & 0.529 \\
uci\_air\_quality\_1D & 1.032 & 1.091 & -5.39\% & -5.29\% & no & 0.919 \\
uci\_air\_quality\_1H & 0.752 & 0.791 & -4.83\% & -5.13\% & no & 0.482 \\
uk\_covid\_nation\_1D/new & 1.985 & 1.997 & -0.61\% & -0.91\% & no & 0.709 \\
\bottomrule
\end{tabular}

\end{table}

\section{Chronos-2 and cross-learning}
\label{app:chronos}

fev's published Chronos-2 results (fev 0.8.0, commit \texttt{f323f6c}) were produced by fev's script
\texttt{examples/chronos-2/evaluate\_model.py}, which calls \texttt{predict\_fev(task, batch\_size=100, as\_univariate=False,
cross\_learning=True)}. We ran that script on a Colab T4 on five tasks where our Chronos-2 failed the pipeline control, with
cross-learning on and off (Table~\ref{tab:chronos}). With cross-learning on, the script reproduces the published SQL to
$4\cdot10^{-6}$; with it off, it reproduces our values. On \texttt{rohlik\_sales\_1W} our value differs by a further 1.2\% because
the protocol applies Google's last-value imputation to every Tier A model. Changing our input preparation (one series per call
against batches of 100, imputed against raw values, synthetic against real timestamps) changes SQL by less than $10^{-8}$.

\begin{table}[h]
\centering
\caption{Chronos-2 SQL on five fev-bench tasks: fev's script with and without cross-learning, fev's published value, and our Tier A
run.}
\label{tab:chronos}
\small
\begin{tabular}{lrrrr}
\toprule
Task & fev script, cross-learning & published & fev script, no cross-learning & our run \\
\midrule
australian\_tourism & 0.6769555 & 0.6769555 & 0.7396415 & 0.7396415 \\
ecdc\_ili & 2.2713938 & 2.2713937 & 2.4665179 & 2.4665179 \\
hospital & 0.6860471 & 0.6860471 & 0.7013574 & 0.7013574 \\
SZ\_TAXI\_1H & 0.3977234 & 0.3977235 & 0.4166707 & 0.4166707 \\
rohlik\_sales\_1W & 1.2741360 & 1.2741411 & 1.1477089 & 1.1338174 \\
\bottomrule
\end{tabular}
\end{table}

\section{TiRex-2 checkpoints}
\label{app:tirex}

The TiRex-2 authors publish a checkpoint described as trained without the fev-bench datasets (\texttt{NX-AI/TiRex-2-fevbench}).
Its weights are byte-identical to those of \texttt{NX-AI/TiRex-2} (all 283
tensors; same source checkpoint SHA-256 \texttt{184b160f\ldots}; bit-identical forecasts). The same checkpoint was uploaded to both
repositories on 2026-07-02, replacing an earlier TiRex-2 checkpoint. The TiRex-2 report describes one checkpoint per benchmark
without stating which one is published as TiRex-2 \citep{podest2026tirex2}. We therefore dropped the planned TiRex-2 contrast on
Tier A and use the public TiRex-2 checkpoint everywhere else. We report this as a provenance note, without drawing conclusions about
the training data; the evidence is released with the code.

\section{Hardware limits}
\label{app:hardware}

On the 15\,GB Tesla T4, \texttt{tfm3-eval} ran out of memory on two Tier A tasks (\texttt{boomlet\_619}, \texttt{boomlet\_1209}),
deterministically at about 14.3\,GB, because symmetric averaging doubles the batch; the package defaults ran on all tasks after the
GPU cache was released between tasks. TiRex-2 used its non-fused sLSTM kernels (the fused kernels need an Ampere GPU) and is not
run beyond 320 steps.
